\section{The Post-\acr{PC} Era and New Computational Requirements}

The transformation of human labor practices presents qualitatively new challenges for society to address.
Computing technologies, as an indispensable instrument for solving most of these challenges, must also
adapt appropriately to meet the needs of emerging economic and social realities.
Moreover, it appears that changes in computing technologies, like socio-economic transformations, follow
the principles of evolution and co-evolution. This perspective is more than mere metaphor—researchers such as Leo Meyerovich and Ariel Rabkin \cite{meyerovichrabkin} document how programming languages evolve under selective pressure from hardware capabilities, business requirements, and competing technologies, much like biological species in a changing environment.

Indeed, the first computers were programmed in machine code.
Later, for convenience, machine instructions were replaced with mnemonics that represented a bijective correspondence to them.
Machine code and mnemonic languages (assembly) are, in their instruction semantics, much closer to programming
physical memory and computer hardware, making them unsuitable for programming abstract processes.
This inconvenience gave rise to the need for a more abstract \acr{ISA}, enabling
the description of processes of broader scientific interest.

Since the first scientific communities to adopt computing
technologies were mathematicians and engineers, it was natural to expect that the first
programming languages beyond mnemonics would have \acr{ISA}s particularly suited for developing and programming
algorithms and mathematical computations. Thus, the first high-level languages such as \acr{Fortran} (1957) and \acr{Algol} (1958) were created primarily
for building algorithms and mathematical computations. They employed \acr{ISA}s close to mathematical notation,
which through parsing were "translated"—or as the established term goes, "compiled"—to assembly, and from there to machine code.

However, these languages proved less suitable for programming business and production tasks
(writing operating systems, inventory management programs, reading data from multiple inputs, etc.).
The search by programmers for a suitable language during this period did not follow any particular pattern but rather
had a somewhat random character. For example, Pascal was developed by Niklaus Wirth in 1970 to teach students structured programming,
\acr{Ada} was developed by the U.S. Department of Defense in 1980 for critical military and aviation systems with high reliability and safety requirements,
and \acr{BASIC} was developed by John Kemeny and Thomas Kurtz in 1964 at Dartmouth College as a simple language for teaching novice programmers.

The most successful proved to be
the C programming language, created by Dennis Ritchie at Bell Labs in 1972. This language possessed the best
combination of \acr{ISA}s suitable for both algorithmic and hardware design. With it, one could conveniently write both algorithmically
formulated problems and system programming tasks (such as writing an operating system—indeed, the C compiler is written in C itself, the so-called "self-hosting compiler").
Additionally, it was sufficiently efficient in terms of performance. Here again, the compiler played a crucial role,
as its task was to reduce C instructions to assembly instructions.

Thanks to the C language,
many different types of languages were created, such as Prolog (1972), C++ (1985), Perl (1987), Python (1991), Ruby (1995), and many others. Some languages resembled C/C++ in syntax
and functionality, while others sought to differentiate themselves from its semantics.

Particularly noteworthy is the Java language (created by James Gosling at Sun Microsystems in 1995),
which resembles C++ but compiles not to machine code but to bytecode, which is then executed by the \acr{JVM} for each computer architecture, making it maximally portable ("write once, run anywhere"). Additionally, Java does not allow direct memory management (automatic garbage collector).

Other languages diverged from the C paradigm, such as JavaScript (created by Brendan Eich at Netscape in 1995 in 10 days) and Python. They completely eliminated
strict static typing of data and focused on \acr{ISA}s optimized for rapid prototyping, dynamic data processing, and interactive development.

Gradually, languages that do not use strict static data types and have simplified syntax with limited
or no direct machine instructions (high level of abstraction) gained prominence in practice. The nature of software products established the paradigm
of object-oriented programming over others. If the path of programming were linear
or governed by some law, one would expect these trends to solidify over time.
However, precisely the opposite occurred.

After the establishment of web technologies and mobile platforms (the "Post-\acr{PC}" era after 2010), the need for more efficient process execution
brought strictly typed languages back into fashion, with Java and C\# becoming most prevalent.
JavaScript responded to this challenge by becoming a backend language through Node.js technology (created by Ryan Dahl in 2009,
later also through Deno and Bun) and through its "pseudo-typing" in TypeScript (created by Microsoft in 2012, adding static types to JavaScript).

Python responded through the use
of \acr{JIT} compilers like Numba or PyPy, as well as through new compilers like Cython, while also introducing optional typing (type hints) from version 3.5 (2015).

But the story does not end here. Writing increasingly efficient programs that process big data
becomes an ever-growing challenge, and one of the simplest working solutions is parallel processing \cite{pattersonhennessy}.
Initially, parallel processing was reduced to multi-threading, later expanding to
interconnected computing machines (cluster computing, distributed computing), and today the paradigm of heterogeneous programming has emerged,
where program execution is distributed between the \acr{CPU} and the \acr{GPU}.

This led once again to the "reign" of C (and C++), which supports \acr{GPU} programming through technologies like \acr{CUDA}, \acr{OpenCL}, \acr{SYCL}, and protocols for
managing \acr{GPU} memory, with strictly typed data being a mandatory requirement
for executing code on \acr{GPU}s (due to hardware-level optimizations and \acr{SIMD}).

Languages like Python find solutions to this problem by
creating compilers for heterogeneous programs using \acr{LLVM} (such as Numba for \acr{CUDA} kernels), while others like TypeScript use \acr{FFI} protocols to call C/C++ code.
It should be noted that all languages use \acr{FFI} protocols or similar mechanisms, since only C/C++ (and other systems languages) can directly manage \acr{GPU} memory and invoke \acr{CUDA}/\acr{OpenCL}/\acr{Vulkan} \acr{API}s.

The attentive reader will easily notice that in this case as well,
the development of programming languages does not follow some strictly defined linear path but is a series of contingencies,
societal needs, and reserves of variability that languages possess to respond to the demands of
emerging environments. Thus, the trajectory of programming language development has a rather zigzag pattern and
is evolutionary in nature, not deterministic. This dynamic resembles the Red Queen Hypothesis from evolutionary biology—programming languages must constantly
adapt not to displace a given paradigm or technology, but to remain competitive in a rapidly changing technological environment. Conversely,
the adaptation of individual languages and computing technologies changes the technological environment itself.

Thus, already-created programming languages do not disappear but, through the forces of interdependence and functional complementarity, specialize in different niches,
such as scientific (\acr{Fortran}, \acr{MATLAB}, Wolfram Language, R), \acr{DSL}
(\acr{SQL} for databases, Verilog/\acr{VHDL} for hardware design, shader languages for \acr{GPU} graphics), and legacy niches
(\acr{COBOL} for financial infrastructure, \acr{Ada} for aviation safety, \acr{Fortran} for climate models). The factor that
determines how a programming language or computing technology will situate itself relative to others is the capacity for maximal adaptation
to the environment and the specialization niche.

Another phenomenon worthy of note is that once again there arises
the need for compilers and compilation infrastructure (such as \acr{LLVM}) to optimize code efficiency for different architectures (\acr{x86}, \acr{ARM}, \acr{RISC-V}, \acr{GPU}). \acr{LLVM} serves as a "universal adapter" in this ecosystem, allowing different languages to compile efficiently for multiple hardware platforms—yet another example of technological co-evolution.

This understanding of computing technologies and interpretation of events in the programming sphere
I largely obtained from the prominent scientist, aristocrat, and public benefactor Peter Sveshtarov through the numerous conversations we have had.
Thanks to him, I was inspired to build this project.

